\capitulo{4}{Metodología}

\section{Marco metodológico y uso de herramientas asistidas}
El desarrollo de este trabajo se ha articulado bajo una metodología de diseño propio centrada en la resolución de problemas clínicos mediante ingeniería de software. Al tratarse del desarrollo de un prototipo funcional se priorizó un enfoque ágil que permitiera validar de forma rápida la lógica clínica del recomendador.

Para optimizar ciertas fases técnicas y garantizar el cumplimiento de estándares de seguridad y legalidad, se ha hecho uso de herramientas de Inteligencia Artificial Generativa (IAG) como soporte técnico especializado. Es importante destacar que la IA ha actuado como un recurso metodológico y consultivo, siendo la toma de decisiones, la validación clínica y la arquitectura final responsabilidad íntegra del autor. El detalle pormenorizado de cada interacción, prompts utilizados y criterios de validación se encuentra documentado en el \textbf{Apéndice I: Registro de interacciones con sistemas de IA}.

\section{Descripción de los datos}
Los datos que sustentan el sistema de recomendación provienen de un análisis de soluciones tecnológicas actuales para la comunicación en pacientes con Esclerosis Lateral Amiotrófica (ELA). La base de conocimiento se ha estructurado en una base de datos relacional que categoriza los Sistemas Aumentativos y Alternativos de Comunicación (SAAC) según criterios clínicos y técnicos.

\begin{itemize}
    	\item \textbf{Origen de los datos:} La información ha sido recopilada a partir de catálogos de productos de apoyo (CEAPAT) \cite{ceapat:productos}, documentación de asociaciones de ELA (como ELA Andalucía) \cite{elaandalucia:guia} y especificaciones técnicas de desarrolladores 		de software especializado (Tobii Dynavox, 	Smartbox, Qinera, Irisbond, etc.) \cite{tobiidynavox:products,qinera:caa,irisbond:soluciones}.
    	\item \textbf{Tipología y formato:} Los datos se almacenan en un sistema de gestión de bases de datos relacionales (PostgreSQL). Se dividen en:
    \begin{itemize}
	\item \textbf{Entidades maestras:} Catálogos de idiomas (español, gallego, catalán, euskera), tipos de entrada (manos, ojos, cabeza, voz, pulsador), plataformas (Windows, iOS, Android, web, panel físico), entornos de uso (domicilio o exterior) y métodos de comunicación (alfabeto o 		pictogramas).
	\item \textbf{Sistemas SAAC:} 46 registros que incluyen software, hardware y servicios de Voice Banking. La selección inicial se basó en los catálogos de CEAPAT y entidades de ELA citadas anteriormente. Esta base fue posteriormente validada y contrastada mediante un análisis 			asistido por IA para asegurar que no se omitieran sistemas críticos en la progresión de la enfermedad, garantizando así un catálogo exhaustivo (ver proceso de validación en \textit{Apéndice I, Interacción 1}). Todas las fuentes técnicas finales han sido referenciadas en la 				bibliografía general.
        	\item \textbf{Requisitos funcionales:} Atributos numéricos que definen los niveles mínimos necesarios (0-3) de capacidad visual, auditiva, cognitiva, tecnológica y de habla para cada sistema.
    \end{itemize}
    	\item \textbf{Carga e inicialización de datos:} Para esta primera versión del prototipo, se ha optado por una estrategia de inicialización estática mediante scripts SQL estructurados para el poblamiento del catálogo de 46 sistemas. Si bien esto asegura un entorno controlado y 			reproducible para la validación del MVP, la automatización de este flujo mediante procesos ETL (Extract, Transform, Load) y el control de versiones de datos se contemplan como evolución del sistema.
	\item \textbf{Características principales:} Los datos permiten un filtrado dinámico mediante consultas SQL que cruzan las capacidades del usuario con las prestaciones del sistema.
\end{itemize}

Una visión más detallada de los datos se puede encontrar en el anexo F-Datos.

\section{Técnicas y herramientas}
El desarrollo de la aplicación se ha basado en una arquitectura web con procesamiento de reglas de recomendación.

\subsection{Herramientas de desarrollo}
\begin{itemize}
    	\item \textbf{Python y Flask:} Seleccionados por la versatilidad del lenguaje y la familiarización adquirida durante el grado. Para robustecer el despliegue frente a vulnerabilidades, la implementación de seguridad (basada en estándares OWASP) contó con el soporte técnico 				documentado en el \textit{Apéndice I (Interacción 3)}, utilizando la librería Flask-Talisman.
    	\item \textbf{Frontend (HTML/CSS):} Se ha diseñado una interfaz mediante pestañas para facilitar la navegación. Mientras que la estructura funcional es de creación propia, el diseño estético y la optimización de la hoja de estilos CSS se apoyaron en herramientas de IA para 				garantizar una visualización profesional y accesible (ver \textit{Apéndice I, Interacción 2}).
    	\item \textbf{PostgreSQL y SQLAlchemy:} Para la persistencia de datos se utiliza PostgreSQL. La interacción entre la aplicación y la base de datos se realiza mediante el ORM SQLAlchemy, abstrayendo las consultas SQL directas y permitiendo una gestión de modelos orientada a 			objetos más intuitiva y segura frente a ataques de inyección SQL.
    	\item \textbf{Jinja2:} Como motor de plantillas para la generación dinámica de las interfaces de usuario (cuestionario e informes de resultados). Es utilizado como extensión al uso de Flask.
\end{itemize}

\section{Metodología software}
Tal y como se introdujo en la sección anterior, el software desarrollado se categoriza como un Sistema de Apoyo a la Decisión (SAD). Su construcción se ha articulado mediante una metodología iterativa e incremental, centrada en la implementación de un motor de decisión basado en lógica proposicional.

\subsection{Modelado de la lógica proposicional}
El recomendador es un sistema que evalúa un conjunto de sentencias lógicas, el conocimiento clínico se traduce en operaciones booleanas y comparativas:

\begin{itemize}
    	\item \textbf{Estructura de las reglas:} Se han definido reglas bajo la estructura condicional IF-THEN. Por ejemplo, si la proposición indica que si el control manual es nulo, se excluyen automáticamente los sistemas de interacción táctil.
    	\item \textbf{Inferencia por filtrado progresivo:} El sistema opera mediante un proceso de eliminación. A medida que el usuario responde al cuestionario, se van validando o invalidando proposiciones lógicas que reducen el conjunto de soluciones posibles hasta obtener el 				subconjunto óptimo para el paciente.
    	\item \textbf{Lógica de post-procesamiento:} Además de la exclusión, se han implementado reglas de recomendación proactiva para sugerir acciones preventivas como el Voice Banking.
\end{itemize}

\subsection{Organización de fases y tareas}
Para garantizar la fiabilidad del sistema, el desarrollo se organizó en las siguientes fases:

\begin{enumerate}
    	\item \textbf{Fase de análisis y captura de reglas:} Identificación de los umbrales críticos de capacidad (visual, cognitiva, motora). Se incluye aquí el diseño de la base legal del sistema. Dada la sensibilidad de los datos de salud, se utilizó IA para la redacción técnica del borrador de la 	cláusula de consentimiento según el RGPD, la cual fue posteriormente revisada y adaptada (ver \textit{Apéndice I, Interacción 5}). Asimismo, las auditorías de seguridad y robustecimiento del sistema (como la integración de Flask-Talisman) se planificaron para mitigar los riesgos 			detectados antes del despliegue final.
    	\item \textbf{Diseño e implementación del motor lógico:} Traducción de las reglas a consultas SQL y estructuras de control en Python (Flask). En esta fase se realizaron múltiples iteraciones para ajustar los filtros lógicos a medida que se realizaban pruebas de funcionamiento.
    	\item \textbf{Validación del sistema experto:} Verificación del comportamiento del recomendador mediante la introducción de perfiles de usuario tipo para asegurar que los resultados fueran clínicamente coherentes. Para esta fase, se generó un banco de pruebas de estrés 				mediante IA para simular perfiles de usuario complejos, asegurando la coherencia clínica de los resultados (ver \textit{Apéndice I, Interacción 4}).
\end{enumerate}


